
Para comprender el significado del coeficiente de dilatación lineal, llene los espacios vacíos que aparecen en las afirmaciones siguientes: Cuando se dice que el coeficiente de dilatación lineal del plomo vale $29 \cdot 10^{-6}$ $^{o}$C$^{-1}$, esto significa que una barra de plomo
\begin{enumerate}
    \item De 1 km de longitud se dilata $29 \cdot 10^{-6}$ km cuando su temperatura aumenta en \rule{60mm}{0.1mm}.
    \item De 1 pulgada de largo se dilata $29 \cdot 10^{-6}$ pulgadas cuando su temperatura aumenta en \rule{60mm}{0.1mm}.
    \item De 1 cm de longitud se dilata \rule{20mm}{0.1mm} cm cuando su temperatura aumenta en 1 $^o$C.
\end{enumerate}
